\documentclass[11pt,a4paper]{article}

\usepackage[utf8]{inputenc}
\usepackage[T1]{fontenc}
\usepackage[spanish,es-noshorthands]{babel}
\usepackage{lmodern}
\usepackage[landscape,margin=1.5cm]{geometry}
\usepackage{xcolor}
\usepackage{tikz}
\usepackage[hidelinks]{hyperref}
\usetikzlibrary{arrows.meta,shapes.geometric}

\definecolor{acento}{RGB}{176,0,32}   % rojo URJC
\definecolor{gris}{RGB}{90,90,90}

\pagestyle{empty}

\tikzset{
  concepto/.style={ellipse, draw=acento, thick, fill=acento!6, align=center,
    inner sep=1pt, minimum height=1.1cm, font=\normalsize},
  central/.style={concepto, fill=acento, text=white, font=\Large\bfseries,
    minimum height=1.8cm},
  linea/.style={thick, gris},
  flecha/.style={linea, -{Stealth[length=2.5mm]}},
  % Verbo compartido: una línea llega al verbo y de él se reparten las flechas
  nexo/.style={fill=white, inner sep=1.5pt, font=\small\itshape, text=acento!80!black},
  verbo/.style={pos=0.55, fill=white, inner sep=1.5pt, font=\small\itshape,
    text=acento!80!black, align=center},
}

\begin{document}

\begin{center}
  {\Large\bfseries\color{acento} Medir lo que no se puede alcanzar: ¿dónde caben más paneles solares en nuestro instituto?}\\[4pt]
  {\small\color{gris} Situación de aprendizaje de Matemáticas I \textperiodcentered{} Diseño Curricular de las Matemáticas \textperiodcentered{} Javier Gómez López}
\end{center}

\vfill
\begin{center}
\resizebox{\linewidth}{!}{%
\begin{tikzpicture}[x=1cm, y=0.85cm]
  % Mapa radial: el reto en el centro y las ramas repartidas a ambos lados.

  % ---------- Concepto central y origen ----------
  \node[central]  (C)   at (0,0)    {Reto de\\los paneles};
  \node[concepto] (dir) at (0,-6.5) {Dirección\\del centro};

  % ---------- Primer nivel ----------
  \node[concepto] (sab) at (-6,5)    {Saberes de\\Matemáticas I};
  \node[concepto] (met) at (-6,-4)   {Metodología};
  \node[concepto] (mod) at (6,5)     {Modelo\\de sombras};
  \node[concepto] (med) at (6,0.5)   {Medición\\en el centro};
  \node[concepto] (eva) at (6,-4.4)  {Evaluación};
  \node[concepto] (div) at (6,-7.5)  {Diversidad\\(DUA)};

  % ---------- Izquierda ----------
  \node[concepto] (vec)  at (-12.5,8.6) {Vectores};
  \node[concepto] (trig) at (-12.5,6.2) {Trigonometría};
  \node[concepto] (tri)  at (-12.5,3.8) {Resolución\\de triángulos};
  \node[concepto] (rec)  at (-12.5,1.4) {Rectas};
  \node[concepto] (mag)  at (-12.5,-1)  {Sesiones\\magistrales};
  \node[concepto] (fich) at (-12.5,-3)  {Talleres\\de GeoGebra};
  \node[concepto] (rol)  at (-12.5,-5)  {Equipos con\\roles rotatorios};
  \node[concepto] (etw)  at (-12.5,-7)  {Proyecto\\eTwinning};

  % ---------- Derecha ----------
  \node[concepto] (prop)  at (12.5,5)    {Propuesta\\de ampliación};
  \node[concepto] (ods)   at (12.5,8.8)  {ODS 7 y 11};
  \node[concepto] (libro) at (12.5,1.2)  {Libro de\\GeoGebra};
  \node[concepto] (exam)  at (12.5,-1.4) {Examen\\individual};
  \node[concepto] (rub)   at (12.5,-3.4) {Libro y\\exposición};
  \node[concepto] (fci)   at (12.5,-5.4) {Fichas y\\cuestionarios};
  \node[concepto] (obs)   at (12.5,-7.4) {Observación y\\co/autoevaluación};

  % ---------- Enlaces ----------
  \draw[flecha] (dir) -- node[verbo] {encarga} (C);

  \draw[flecha] (C.north west) to[out=125, in=-55] node[verbo] {moviliza} (sab);
  \draw[flecha] (C) to[out=10, in=180] node[verbo] {exige} (med);
  \draw[flecha] (C) to[out=-145, in=0] node[verbo] {se apoya en} (met);
  \draw[flecha] (C) to[out=-20, in=180] node[verbo] {se comprueba con} (eva);
  \draw[flecha] (C) to[out=-55, in=180] node[verbo, pos=0.7] {atiende la} (div);

  \node[nexo] (ab) at (-9.3,5) {abarcan};
  \draw[linea]  (sab) -- (ab);
  \draw[flecha] (ab.north) to[out=100,  in=0] (vec);
  \draw[flecha] (ab.west)  to[out=155,  in=0] (trig);
  \draw[flecha] (ab.west)  to[out=-155, in=0] (tri);
  \draw[flecha] (ab.south) to[out=-100, in=0] (rec);

  \draw[flecha] (sab.north east) to[out=40, in=140] node[verbo, pos=0.5] {construyen} (mod.north west);
  \draw[flecha] (med) -- node[verbo] {aporta datos al} (mod);
  \draw[flecha] (mod)  -- node[verbo] {justifica} (prop);
  \draw[flecha] (prop) -- node[verbo] {contribuye a} (ods);
  \draw[flecha] (prop) -- node[verbo] {se recoge en} (libro);

  \node[nexo] (comb) at (-9.3,-4) {combina};
  \draw[linea]  (met) -- (comb);
  \draw[flecha] (comb.north) to[out=100,  in=0] (mag);
  \draw[flecha] (comb.west)  to[out=155,  in=0] (fich);
  \draw[flecha] (comb.west)  to[out=-155, in=0] (rol);
  \draw[flecha] (comb.south) to[out=-100, in=0] (etw);

  \node[nexo] (pon) at (9.3,-4.4) {incluye};
  \draw[linea]  (eva) -- (pon);
  \draw[flecha] (pon.north) to[out=80,  in=180] (exam);
  \draw[flecha] (pon.east)  to[out=25,  in=180] (rub);
  \draw[flecha] (pon.east)  to[out=-25, in=180] (fci);
  \draw[flecha] (pon.south) to[out=-80, in=180] (obs);
\end{tikzpicture}%
}
\end{center}
\vfill

{\footnotesize\color{gris} Fuente: elaboración propia a partir de la situación de aprendizaje «Medir lo que no se puede alcanzar» (1.º de Bachillerato, Matemáticas I, 20 sesiones). IES Ejemplo es un centro ficticio.}

\end{document}
